\documentclass[10pt,a4paper]{amsart}
\usepackage[margin=19mm]{geometry}
\usepackage{amsmath,amssymb,mathtools}
\usepackage[scr=rsfs,cal=euler]{mathalfa}
\usepackage{xfrac}
\usepackage[unicode,hidelinks,bookmarks=false]{hyperref}
\newcommand{\ie}{i.e.\ }
\newcommand{\cf}{cf.\ }
\newcommand{\resp}{respectively\ }
\newcommand{\Subsection}{\subsection}
\providecommand{\nobreakdash}{\nobreak\hbox{-}}
\setlength{\emergencystretch}{2em}
\multlinegap0pt
\usepackage{bm}
\usepackage{bbm}
\usepackage{dsfont}
\usepackage{xspace}
\usepackage{cancel}
\usepackage{mathtools}
\usepackage{amsthm}
\makeatletter
\@ifundefined{theorem}{\newtheorem{theorem}{Theorem}}{}
\@ifundefined{proposition}{\newtheorem{proposition}[theorem]{Proposition}}{}
\makeatother
\newcommand{\R}{{\ensuremath{\mathbb{R}}}}
\title[HOT-POT: Optimal Transport for Sparse Stereo Matching]{HOT-POT: Optimal Transport for Sparse Stereo Matching}
\author{Antonin Clerc, Michael Quellmalz, Moritz Piening, Philipp Flotho, Gregor Kornhardt, Gabriele Steidl}
\date{Source: arXiv:2601.12423v2 (CC BY 4.0)}
\begin{document}
\maketitle

\noindent\emph{Source and licence.} HOT-POT: Optimal Transport for Sparse Stereo Matching. Version arXiv:2601.12423v2 (CC BY 4.0), distributed under the Creative Commons Attribution 4.0 International licence, as recorded in the arXiv licence metadata for this exact version and shown on the abstract page https://arxiv.org/abs/2601.12423v2. Full rights notice: licenses/arxiv-2601.12423-CC-BY-4.0.txt.\par\smallskip

\noindent\textit{Editorial scope.} Counted unit: the Proposition whose source label is \texttt{prop:ray\_distance\_pseudo\_definite\_property} in the article (source0.tex lines 552--559, source label \texttt{prop:ray\_distance\_pseudo\_definite\_property}), delivered together with the article's own complete original proof of it (source0.tex lines 560--575, one \texttt{proof} environment of 16 lines). No mathematics is written, repaired or re-derived: statement and proof are the article's own text line for line. Required context retained uncounted: eq:model (lines 418--422); eq:front (lines 426--428); eq:lines (lines 471--475). Every definition, display and label of those lines is retained unchanged. Omitted from this package and not redistributed: 1--417, 423--425, 429--470, 476--551, 576--1564. Nothing inside a retained excerpt is omitted or altered: every retained body is delivered exactly as the article's lines are written, so no omission receipt of any kind accompanies this package. Citations: the retained excerpts carry no citation command at all, so no bibliography entry of the article is transcribed and no reference is redistributed. Reference adaptation: none was needed and none was made; every reference inside the delivered unit resolves inside it, so no unresolved reference or citation is produced. Printed slips kept literal: no printed word, symbol, hypothesis or number of the counted statement or of its proof was altered. Layout and class: the article's own class is \texttt{article} with packages including geometry, lmodern, fontenc, inputenc, microtype, titlesec, etoolbox, booktabs, amsmath, amsfonts, amssymb, amsthm, mathtools, bm; the wrapper is the lane's pinned \texttt{amsart} editorial wrapper at 10pt on A4, which restates the theorem-like environments the retained text needs and adds the article's own macro definitions that the retained text needs (\texttt{\textbackslash R}), with extra wrapper packages (bm, bbm, dsfont, xspace, cancel, mathtools). The delivered numbering is the unit's own. Assets: no retained span contains a figure environment, a graphics-inclusion command, a PDF-page inclusion, a table or a TikZ picture; no image file is copied into this package, the package declares no asset dependency and no image file is redistributed with it. Licence: version arXiv:2601.12423v2 of this article is distributed under the Creative Commons Attribution 4.0 International licence, as recorded in the arXiv licence metadata for this exact version and shown on the abstract page https://arxiv.org/abs/2601.12423v2. A full rights notice is delivered at licenses/arxiv-2601.12423-CC-BY-4.0.txt. Characters: every retained excerpt is pure ASCII, so the delivered engine, XeLaTeX with its default Latin Modern text font, reports no missing character.
\begin{equation} \label{eq:model}
\lambda_l x = K_l w 
\quad \text{and} \quad
\lambda_r y = K_r (Rw + t),
\end{equation}

\begin{equation} \label{eq:front}
\langle w, e_3\rangle>0 \quad \text{and} \quad \langle Rw+t,e_3\rangle>0.
\end{equation}

\begin{equation} \label{eq:lines}
        w = \lambda_l \tilde x 
        \quad \text{and} \quad
        w = \lambda_r R^\top \tilde y - R^\top t.
\end{equation} 

\begin{proposition}
\label{prop:ray_distance_pseudo_definite_property}
    If $w\in\mathcal W$ and $x,y\in\mathbb P^2$ are given by \eqref{eq:model},
    then $d^\mathrm{ray}(x,y)=0$.
    Conversely, if $x,y\in\mathbb P^2$ and $d^\mathrm{ray}(x,y)=0$, 
    there exists $w\in\R^3$ such that \eqref{eq:model} holds for some $\lambda_l,\lambda_r\in\R$. 
    If additionally $r_x \times r_y \neq 0$, then $w$ is uniquely determined.
\end{proposition}

\begin{proof}
Let $w$ satisfy \eqref{eq:front} and $x,y$ be the projections of $w$ via \eqref{eq:model}. 
By construction, 
$w$ is on both lines $\mathcal L_x$ and $\mathcal L_y$.
If the lines are not identical, i.e., if $r_x\times r_y\neq0$, we have $w=b_x=b_y$.
By \eqref{eq:front}, we are in the first case of the definition of $d^\mathrm{ray}$ and hence we have $d^\mathrm{ray}(x,y)=0$.
Otherwise, if the lines are identical, i.e., if $r_x\times r_y=0$, their distance is zero and so $d^\mathrm{ray}(x,y)=0$.

Conversely, let $d^\mathrm{ray}(x,y)=0$.
We note that $\|t\|$ does not vanish by assumption.
If $0=r_1\times r_2 = \tilde x\times R^\top\tilde y$, we have $\tilde x\times R^\top t=0$.
Hence, $\tilde x$, $R^\top\tilde y$, and $R^\top t$ are all located one line through the origin and neither does vanish.
Hence, \eqref{eq:lines}, which is equivalent to \eqref{eq:model}, is fulfilled for $w=\alpha \tilde x$ with any $\alpha\neq0$ and some $\lambda_l,\lambda_r\in\R$.
Otherwise, the lines intersect in one point $b_x=b_y$, and hence \eqref{eq:lines} is fulfilled with $w=b_x=b_y$.
In this case, $w$ satisfies~\eqref{eq:front}.
\end{proof}

\end{document}
